\documentclass[a4paper]{article}

\usepackage{graphicx}
\usepackage{amsmath,amssymb}
\usepackage{minted}
\usepackage{multirow}
\usepackage{float}
\usepackage{todonotes}
\usepackage{forest}
\usepackage{xstring}
\usepackage{xcolor}
\usepackage[most]{tcolorbox}
\usepackage{xparse}
\usepackage{natbib}

\usetikzlibrary{graphs,graphdrawing}
\usegdlibrary{layered}

% for external graphviz
\input{/home/donkarlo/Dropbox/repo/nd_language_ai_project/src/nd_language_ai/formal/computer/programming/tex/latex/code_clips/external_graphviz.tex}

% for tags
\input{/home/donkarlo/Dropbox/repo/nd_language_ai_project/src/nd_language_ai/formal/computer/programming/tex/latex/parts/control_sequence/kind/semtag/semtag.tex}

% for links
\input{/home/donkarlo/Dropbox/repo/nd_language_ai_project/src/nd_language_ai/formal/computer/programming/tex/latex/code_clips/seehere.tex}

\title{Bayesian sensory interation}
\date{}
\author{Mohammad Rahmani}

\begin{document}
    \maketitle
    This study investigates how visual optic-flow and vestibular signals are combined when humans estimate their direction of self-motion. Eight participants completed heading-discrimination tasks under visual-only, vestibular-only, and combined conditions containing spatial conflicts of $\pm 6^{\circ}$ or $\pm 10^{\circ}$.
    Combining the two sensory cues reduced perceptual variance approximately as predicted by maximum-likelihood estimation, even when the cues indicated conflicting heading directions. However, the observed cue weights were not determined solely by the reliabilities measured in the unimodal conditions, because participants consistently gave greater weight to vestibular information. The results suggest that visual--vestibular integration is statistically near-optimal in terms of precision, but that its systematic vestibular bias requires a Bayesian prior or a more comprehensive model to explain \cite{butler-2010-bayesian-integration-of-visual-and-vestibular-signals-for-heading}.
    \bibliographystyle{plainnat}
    \bibliography{/home/donkarlo/Dropbox/repo/nd_shared_research/bibliography/bibliography.bib}
\end{document}